\section{结论与后续工作}

\subsection{综合结论}

\texttt{YH\_rv\_cpu} 已经形成较完整的初赛工程形态：架构层面采用边界清晰的五级顺序流水，
验证层面建立了从 smoke、baseline、扩展覆盖到 CoreMark、\texttt{impl50} 和 FPGA-like probe 的证据链，
原型层面完成了可复现的 FPGA 原型适配。综合来看，本设计在初赛阶段具备三项较明确的特点：
\begin{itemize}
    \item 方案结构清晰，能够较好地说明设计原理与工程取舍；
    \item 验证矩阵完整，实验结果具备较强的可复核性；
    \item 资源与时序结果能够支撑后续 FPGA 原型落地，而不是停留在概念设计阶段。
\end{itemize}

\subsection{后续工作}

后续工作主要包括：
\begin{itemize}
    \item 在实体板卡到位后完成 UART/LED 连线、boot banner 观测和最终板级 I/O 约束确认；
    \item 在保持当前验证基线稳定的前提下，继续开展控制流 redirect 成本分析与性能优化实验；
    \item 结合后续赛程需要，评估是否引入更多扩展指令或进一步增强原型系统能力。
\end{itemize}

\subsection{结语}

总体而言，\texttt{YH\_rv\_cpu} 的价值不在于追求极端复杂结构，
而在于以清晰的设计边界、完整的验证矩阵和可复现的实现结果，建立了一套具备工程完整性的处理器方案。
